% =====================================================================
%  Hiện thực UC07, UC07, UC07 (chèn vào hien-thuc-uc.tex sau mục UC06)
% =====================================================================
\subsection{Xóa thành viên, chuyển quyền và rời nhóm (UC07, UC07, UC07)}\label{sssec:ht-roi-nhom}

Ba use case do \texttt{household.Service} hiện thực qua các API dưới \path{/spaces/{sid}}: G08 \texttt{DELETE} \path{.../members/{userId}} (\texttt{RemoveMember}), G09 \texttt{POST} \path{.../leader-transfer} (\texttt{TransferLeader}) và G10 \texttt{POST} \path{.../leave} (\texttt{LeaveSpace}). Hình~\ref{fig:seq-roi-nhom-chuyen-quyen} theo con đường dài nhất của \uc{UC07}: trưởng nhóm rời nhóm còn thành viên bị từ chối (bước C2a), chuyển quyền theo \uc{UC07}, rồi rời nhóm với tư cách thành viên; G08 đi theo nhánh rời nhóm khi còn thành viên khác.

\diagram[\textwidth]{c3-36-seq-roi-nhom-chuyen-quyen}{Biểu đồ tuần tự hiện thực use case Chuyển quyền trưởng nhóm, Rời nhóm và Xóa thành viên}{fig:seq-roi-nhom-chuyen-quyen}

Cả ba API lấy khóa theo thứ tự toàn cục của mục~\ref{sssec:kt-fefo}, giống G07 ở mục~\ref{sssec:ht-tham-gia-nhom}: khóa lũy đẳng, khóa tư vấn của không gian (\texttt{pg\_advisory\_xact\_lock} trên \texttt{hashtext(space\_id::text)}), các dòng \texttt{space\_members} theo \texttt{user\_id} tăng dần, các mặt hàng theo \texttt{id}, sau cùng \texttt{spaces.change\_seq}. Khóa tư vấn tuần tự hóa mọi thay đổi thành viên, vai trò và mã mời của một nhóm, nên hai thiết bị của trưởng nhóm chuyển quyền cho hai người khác nhau, hay chuyển quyền cho người đang rời nhóm, không thể đan xen. Vai trò người gọi được đọc lại trong giao dịch thay vì lấy từ bộ đệm \texttt{member:\{sid\}:\{uid\}} có thể chậm 60 giây; người không còn là trưởng nhóm nhận 403 mã 3007 kèm \msg{MSG13} (\uc{UC07} bước 4a). Khóa chết hoặc chờ khóa quá hạn được thử lại tối đa ba lần rồi trả 503 mã 5003.

G09 khóa hai dòng của người chuyển và người nhận; người nhận vừa rời nhóm (\uc{UC07} bước B4a) làm yêu cầu nhận 422 mã 4004 kèm \br{BR05} để ứng dụng tải lại danh sách. Thứ tự hai câu cập nhật là bắt buộc: hạ người chuyển xuống \texttt{member} trước, nâng người nhận lên \texttt{leader} sau (bước 5). \texttt{uq\_space\_leader}, tức \texttt{UNIQUE(space\_id) WHERE role = 'leader'}, là chỉ mục duy nhất từng phần nên được kiểm tra sau từng câu lệnh, không trì hoãn được như ràng buộc \texttt{DEFERRABLE}~\cite{postgresql}; nâng trước sẽ tạo hai trưởng nhóm và bị từ chối. Trạng thái trung gian không có trưởng nhóm chỉ tồn tại trong giao dịch, nên bên ngoài luôn thấy đúng một trưởng nhóm (\br{BR05}); nếu một luồng mã bỏ qua khóa tư vấn, lỗi 23505 trên chỉ mục này trả 409 mã 4002 kèm \msg{MSG35}. Cùng giao dịch thu hồi mã mời do trưởng nhóm cũ tạo (\br{BR06}) và chèn \texttt{PushNotification} cho \texttt{leader.transferred} (bước 6); sau commit, khóa đệm vai trò của hai người bị xóa và \texttt{space:\{sid\}} được phát, nên giao diện của trưởng nhóm cũ mất các nút quản lý trong khoảng 3 giây (\nfr{NFR10}).

G10 kiểm tra \br{BR07} sau khi khóa: trưởng nhóm của nhóm còn thành viên khác nhận 422 mã 4004 kèm \msg{MSG24} (bước 12), ứng dụng mở \uc{UC07} rồi gửi lại G10. Thành viên duy nhất (bước 2b), sau cảnh báo \msg{MSG19} trên ứng dụng, làm giao dịch xóa dòng thành viên, chèn tác vụ \texttt{PurgeSpace} và đặt \texttt{dissolved\_at} ngay trong câu tăng \texttt{change\_seq} cuối giao dịch (bước 13). Mọi truy vấn lọc \texttt{dissolved\_at IS NULL} nên nhóm biến mất ngay, còn \texttt{PurgeSpace} lũy đẳng xóa mềm dữ liệu chung và tài liệu Elasticsearch của không gian ở phía sau (\nfr{NFR18}). Khi nhóm còn người khác, dịch vụ xóa dòng thành viên rồi gọi \texttt{household/api.MemberLeaveHook} do shopping hiện thực trong cùng giao dịch (bước 14--16): mặt hàng \texttt{assigned} của người rời được khóa theo \texttt{id}, về \texttt{to\_buy}, bỏ \texttt{assignee\_id}, ghi \texttt{space\_changes} và phát \texttt{item.unassigned} (\br{BR07}). Xóa dòng thành viên trước là có chủ đích: phân công S10 đọc dòng người được giao bằng \texttt{FOR SHARE} (Hình~\ref{fig:seq-phan-cong}), nên giao dịch rời nhóm chờ phân công đang chạy kết thúc rồi hoàn trả cả mặt hàng vừa giao, còn phân công đến sau không thấy người được giao và nhận 422. Giao dịch cũng chuyển \texttt{last\_space\_id} của người rời về không gian cá nhân qua cổng của identity và chèn \texttt{PushNotification} cho \texttt{member.left} để báo trưởng nhóm (bước 17). Sau commit, dịch vụ phát thêm \texttt{user:\{uid\}} để bản API giữ WebSocket của người rời hủy đăng ký kênh của nhóm (\br{BR23}). Dữ liệu người rời đã đóng góp vẫn thuộc nhóm vì gắn với \texttt{space\_id}. G08 dùng đúng nhánh này, khác ở chỗ người gọi phải là trưởng nhóm, không được tự xóa mình (422 mã 4004) và sự kiện \texttt{member.removed} báo cho người bị xóa. G08, G10 trên không gian cá nhân trả 422 mã 4004 (\br{BR04}).
